\documentclass[a4paper,10pt]{report}\input{../0.Préambule/Préambule_cours_prof}\begin{document}

\newpage
\section*{SCRATCH - Exemples de scripts de tracés}
\addcontentsline{toc}{section}{SCRATCH - Exemples de scripts de tracés}

\subsection*{Initialisation d'un tracé}
\addcontentsline{toc}{subsection}{Initialisation d'un tracé}

\vspace{-5mm}
\noindent \begin{minipage}{11cm}
Il est important de bien se positionner avant de faire un tracé. On peut :
\begin{enumerate}
\item[•] réduire la taille du lutin afin de mieux voir le tracé
\item[•] se placer au centre du cadre (ou ailleurs)
\item[•] regarder vers la droite
\item[•] tout effacer
\item[•] et se mettre en position pour tracer
\end{enumerate}
\end{minipage}
\hspace{5mm}
\begin{minipage}{6cm}
\setscratch{scale=0.8}
\begin{scratch}
\blockinit{quand \greenflag est cliqué}
\blocklook{mettre à \ovalnum{20} $\%$ de la taille initiale}
\blockmove{aller à x: \ovalnum{0} y: \ovalnum{0}}
\blockmove{s'orienter à \ovalnum{90}}
\blockpen{effacer tout}
\blockpen{stylo en position d'écriture}
\end{scratch}
\end{minipage}

\vspace{10mm}
\noindent\begin{minipage}[t]{10.8cm}
\subsection*{Dimensions de la scène}
\addcontentsline{toc}{subsection}{Dimensions de la scène}

\includegraphics[scale=0.315]{../40.Image/SCRATCH_cours_2_1.png}
\end{minipage}
\begin{minipage}[t]{5.6cm}
\subsection*{Tracé d'un rectangle}
\addcontentsline{toc}{subsection}{Tracé d'un rectangle}

\setscratch{scale=0.8}
\begin{scratch}
\blockrepeat{répéter \ovalnum{2} fois}
{
	\blockmove{avancer de \ovalnum{3}}
	\blockmove{tourner de \turnright{} de \ovalnum{90} degrés}
	\blockmove{avancer de \ovalnum{1}}
	\blockmove{tourner de \turnright{} de \ovalnum{90} degrés}	
}
\end{scratch}

\medskip
\noindent\grandscarreaux{\linewidth}{2.4cm}
\end{minipage}

\vspace{10mm}
\noindent\begin{minipage}[t]{8.5cm}
\subsection*{Tracé d'un carré}
\addcontentsline{toc}{subsection}{Tracé d'un carré}

\noindent\begin{minipage}{4cm}
\setscratch{scale=0.8}
\begin{scratch}
\blockrepeat{répéter \ovalnum{4} fois}
{
	\blockmove{avancer de \ovalnum{2}}
	\blockmove{tourner de \turnright{} de \ovalnum{90} degrés}
}
\end{scratch}
\end{minipage}
\begin{minipage}{4cm}
\noindent\grandscarreaux{\linewidth}{3.2cm}
\end{minipage}

\bigskip
\noindent\begin{minipage}{4cm}
\setscratch{scale=0.8}
\begin{scratch}
\blockrepeat{répéter \ovalnum{4} fois}
{
	\blockmove{avancer de \ovalnum{2}}
	\blockmove{tourner de \turnleft{} de \ovalnum{90} degrés}
}
\end{scratch}
\end{minipage}
\begin{minipage}{4cm}
\noindent\grandscarreaux{\linewidth}{3.2cm}
\end{minipage}

\end{minipage}
\begin{minipage}[t]{8.5cm}
\subsection*{Tracé d'un triangle équilatéral}
\addcontentsline{toc}{subsection}{Tracé d'un triangle équilatéral}

\noindent\begin{minipage}{4.2cm}
\setscratch{scale=0.8}
\begin{scratch}
\blockrepeat{répéter \ovalnum{3} fois}
{
	\blockmove{avancer de \ovalnum{3}}
	\blockmove{tourner de \turnright{} de \ovalnum{120} degrés}
}
\end{scratch}
\end{minipage}
\begin{minipage}{4cm}
\noindent\grandscarreaux{\linewidth}{3.2cm}
\end{minipage}

\bigskip
\noindent\begin{minipage}{4.2cm}
\setscratch{scale=0.8}
\begin{scratch}
\blockrepeat{répéter \ovalnum{3} fois}
{
	\blockmove{avancer de \ovalnum{3}}
	\blockmove{tourner de \turnleft{} de \ovalnum{120} degrés}
}
\end{scratch}
\end{minipage}
\begin{minipage}{4cm}
\noindent\grandscarreaux{\linewidth}{3.2cm}
\end{minipage}
\end{minipage}
\end{document}